\begin{frame}
\frametitle{Overview}
\begin{itemize}
  \item Etymology of the term 'Navigation':  Latin $navis$='ship' and $agere$='to move or direct'
  \item Arose out of marine or ship navigation due to the challenge of navigating without landmarks
  \item Modern distinction between navigation and guidance
  \begin{itemize}
    \item navigation associated with state determination (open loop)
    \item guidance associated with control feedback (closed loop)
  \end{itemize} 
\end{itemize} 
\end{frame}

\begin{frame}
  \frametitle{The Navigation Problem from a Mathematical Perspective}
With respect to some mathematical reference frame, determine ...
\begin{itemize}
  \item vector displacement (3d position) and its time derivatives (velocity, acceleration, ...)
  \item rotational orientation (3d attitude) and its time derivatives (angular velocity, angular acceleration, ...) 
  \item Note:  This is a pure mathematics problem ... sensors do not enter into it
\end{itemize} 
\end{frame}


\begin{frame}
  \frametitle{The Navigation Problem from a Mathematical Perspective}
\begin{itemize}
  \item Reference frame is crucial:
  \begin{itemize}
    \item Position:  origin and three linearly independent vectors yields three numbers
    \item Attitude:  rotations about three linearly indepedent vectors yields three numbers
  \item I.e., 6 DOF = six degrees of freedom (3 position coordinates + 3 attitude coordinates)
  \end{itemize}
\end{itemize} 
\end{frame}


\begin{frame}
 \frametitle{The Navigation Problem from a Real World Physical Perspective}
\begin{itemize}
  \item Which reference frame should we use?
  \item Rectilinear or curvilenar coordinate system?
\end{itemize} 
\end{frame}

\begin{frame}
 \frametitle{The Navigation Problem from a Real World Physical Perspective}
\framesubtitle{Position}
For terrestrial navigation, global spherical (spheroidal) coordinates are used:
\begin{itemize}
  \item Latitude:  $\pm 90$ degrees north(+) or south(-) of the equator.
  \item Longitude: $\pm 180$ degrees east(+) or west(-) of the prime meridian through Greenwich, England
  \item Height:  Above the ellipsoid (normal to surface), or radius from center of earth
\end{itemize} 
\end{frame}

\begin{frame}
 \frametitle{The Navigation Problem from a Real World Physical Perspective}
\framesubtitle{Direction:  Bearing, Course, Track}
\begin{itemize}
  \item North Reference:  True and Magnetic. Direction angles are usually referenced clockwise with respect to north.
  \item  Bearing:  The angle with respect to north from one point on earth to another
  \item  Track:    The angle made by velocity vector with respect to north
  \item  Course:   Sometimes used synonym for track; sometimes expresses desired track vs. actual track
  \item  Heading:  Direction vehicle is pointing in with respect to north
\end{itemize} 
\end{frame}

\begin{frame}
 \frametitle{The Navigation Problem from a Real World Physical Perspective}
\framesubtitle{The Need for Sensors}
\begin{itemize}
  \item Reference frames are necessary, but not sufficient for real world navigation
  \item Also need real world sensor measurement data indicative of locations, distances, angles, etc. with respect to these reference frames
  \item But real world sensors are not perfect!  They only provide apparent positions, distances, angles, etc.
\end{itemize} 
\end{frame}
